\documentclass[../../main.tex]{subfiles}
\begin{document}


{\bf [解]}
\paragraph{(1)}

$x^2-3x+5=0$の2解を$\alpha,\beta$とすると解と係数の関係より
\begin{align}
 \begin{cases}
  \alpha+\beta=3 \\
  \alpha\beta=5
 \end{cases} \label{eq:1}
\end{align}
である．新しく
\begin{align}
 \begin{cases}
  b_n=\alpha^n+\beta^n \\
  a_n=b_n-3^n
 \end{cases} \label{eq:2}
\end{align}
とおくと，\cref{eq:1}より$b_n$の漸化式は
\begin{align}
 b_{n+2}
  &=\alpha^{n+2}+\beta^{n+2} \\
  &=(\alpha+\beta)b_{n+1}-\alpha\beta b_n \\
  &=3b_{n+1}-5b_n \label{eq:3}
\end{align}
である．また，$n=1,2$の時
\begin{align}
 \begin{cases}
   b_1 = \alpha + \beta = 3 \\
   b_2 = \alpha^2 + \beta^2 = (\alpha+\beta)^2 - 2\alpha\beta = -1 
 \end{cases} \label{eq:4}
\end{align}
である．
以下$2$項づつの数学的帰納法によりすべての正の整数$n$で$a_n\equiv0\pmod5$であることを示す．

まず$n=1,2$の時，\cref{eq:4}より
\begin{align}
 \begin{cases}
  a_1 = b_1 - 3^{1} = 3-3 = 0 \equiv0\pmod5 \\
  a_2 = b_2 - 3^{2} = -1 -9 = -10 \equiv0\pmod5
 \end{cases}
\end{align}
より成立する．

そこで$n=k,k+1$（$k\in\mathbb N$）での成立を仮定する．
\cref{eq:3}より
\begin{align}
a_{k+2}
 &=b_{k+2}-3^{k+2} \\
 &=3b_{k+1}-5b_k-3^{k+2}\\
 &=3(a_{k+1}+3^{k+1})-5(a_k+3^k)-3^{k+2} \\
 &=3a_{k+1}-5a_k-5\cdot3^k
\end{align}
だが，仮定より$a_{k+1}\equiv a_k\equiv0$で，かつ$-5\cdot3^k$は$5$の倍数だから$a_{k+2}$も$5$の倍数である．よって$n=k+2$でも仮定が成立する．

以上より数学的帰納法によりすべての正の整数$n$で$\alpha^n+\beta^n-3^n$は$5$の倍数である．

\paragraph{(2)}

4種類の目$a,b,c,d$が出る場合，6個の出目の内訳（重複度の分割）は$3+1+1+1$または$2+2+1+1$の2パターンに限る．6個のさいころを区別すれば全事象は$6^6$通りで同様に確からしい．

\subparagraph{$1^\circ$ 内訳が$3,1,1,1$のとき}

$4$個の目のうちどれが$3$回出るかで$6$通り，残り$3$つの目を残り$5$種類の目から選ぶ方法が${}_5C_3=10$通りだから，目の選び方は$6\times10=60$通りある．

各選び方についてどのさいころがどの目になるかの割り当ては，まず$6$つのうちで$3$つある同じ目を選ぶのが${}_6C_{3}$通り．さらに残り$3$つの目の割り当てが$3!$通りで，合計${}_6C_3\cdot3!=20\times6=120$通りである．

以上からこのパターンの場合の数は
\begin{align}
 60\times120=7200=2\cdot6^2\cdot100
\end{align}
である．

\subparagraph{$2^\circ$ 内訳が$2,2,1,1$のとき}

$2$回出る目の選び方が${}_6C_2=15$通り，残りから$1$回ずつ出る目の選び方が${}_4C_2=6$通りで，目の選び方は$15\times6=90$通りある．

割り当てはまず$6$つのうちで$2$つある同じ目を選ぶのが${}_6C_2\cdot{}_4C_2$通りで，残りの$2$つの目を割り当てるのが$2!$通りだから合計で${}_6C_2\cdot{}_4C_2\cdot2!=15\times6\times2=180$通りである．

以上からこのパターンの場合の数は
\begin{align}
 90\times180=16200=2\cdot9^2\cdot100
\end{align}
である．
\casesend

以上より，求める確率は
\begin{align}
\dfrac{2\cdot6^2\cdot100+2\cdot9^2\cdot100}{6^6}=\dfrac{(72+162)\times100}{46656}=\dfrac{23400}{46656}=\dfrac{325}{648}
\end{align}
である．


\newpage

\end{document}
